\chapter{Renormalization, Naturalness, and the Hierarchy Problem}

\lectureoverview{The central ideas of this course grew out of puzzles about
renormalization in the Standard Model. We begin, therefore, with the physical
meaning of renormalization: eliminate degrees of freedom that are too small
and too fast for the question at hand, then encode their effects in an
effective description. A molecule supplies the prototype. Dimensional
analysis then carries the same reasoning into scalar quantum field theory,
where it exposes the quadratic sensitivity of a scalar mass, the hierarchy
problem of the Higgs field, the protection enjoyed by fermions and gauge
bosons, and the still more severe problem of vacuum energy.}

\section{Physics at the Scale of the Question}

Most of the ideas to be developed this quarter originated in questions about
renormalization---sometimes puzzles, sometimes useful observations---in the
Standard Model. Renormalization combines two habits
of thought. The first is to remove from a description structures at distances
so short that they are irrelevant to the question being asked. The second is
dimensional analysis, which often determines the short-distance answer before
any difficult integral is evaluated.

This elimination of detail happens throughout physics. If we wish to
understand a nucleus, we may begin with quarks and quantum chromodynamics.
After calculating the masses and spins of the proton and neutron, together
with their effective nuclear forces, we can forget their quark substructure
for many purposes. Nuclear motion is slow enough that a nonrelativistic
quantum mechanics of protons and neutrons is often the more useful
description.

At the next level, atomic physics does not ordinarily need the internal
constitution of a nucleus. Nuclear physics tells us which nuclei and isotopes
exist and supplies such properties as mass and electric charge. We then begin
again with a new set of effective elementary objects: nuclei and electrons.
Quarks, protons, and neutrons have not ceased to exist; their effects have
been compressed into the parameters carried by the nuclei.

The amount retained depends on the next question. An atomic calculation may
need the nuclear mass because it enters the reduced mass, and it certainly
needs the charge because that determines the Coulomb field. It usually does
not need to know which quark carries which fraction of the nuclear momentum.
There are roughly ninety elemental charges, with more nuclear species once
isotopes are counted. That catalogue can be taken as input. Electrons, on the
other hand, cannot yet be eliminated: their motion is exactly what makes
atomic structure.

After solving the atomic problem, we may similarly regard atoms as the
building blocks of molecules. Susskind's whimsical list of cookie atoms,
coffee atoms, chocolate atoms, and ketchup atoms makes the serious point that
a new description begins at every level. Once molecular properties and
forces are known, molecules can be joined into larger Tinkertoy-like
assemblies without reopening the electron problem at every step.

Each passage is a coarse-graining. The new account may be less exact in an
absolute sense, but it is far more useful at its intended scale. Quantum field
theory makes this issue unavoidable because a field has modes at every
wavelength, including arbitrarily short ones. We therefore sum the effects of
wavelengths too short to resolve and replace them by new effective parameters.
That operation, together with dimensional analysis, is the physical content
of renormalization.

\section{A Molecule as the Prototype}

Consider two atoms and ask for a simple description of their relative
motion. Each atom contains a heavy nucleus and a
cloud of electrons. Even the proton is roughly two thousand times heavier
than an electron, and heavier nuclei increase the separation of scales. A
useful picture is a pair of bowling balls surrounded by rapidly moving flies:
the nuclei respond slowly while the electronic state rearranges quickly.

Let the nuclear positions and momenta be \(R_a\) and \(P_a\), with common
mass \(M\), and let the electron coordinates and momenta be \(r_i\) and
\(q_i\), with mass \(m_e\). Suppressing conventional factors such as
\(4\pi\epsilon_0\), the Hamiltonian has the schematic form
\begin{align}
H={}&
\frac{P_1^2}{2M}+\frac{P_2^2}{2M}
+\frac{e^2}{R_{12}}
+\sum_i\frac{q_i^2}{2m_e}
+\frac12\sum_{i\ne j}\frac{e^2}{r_{ij}}
\nonumber\\
&-\sum_i\left(
\frac{e^2}{|R_1-r_i|}
+\frac{e^2}{|R_2-r_i|}
\right).
\label{eq:molecular-hamiltonian}
\end{align}
The first line contains the nuclear kinetic energies, their mutual Coulomb
repulsion, and the electronic kinetic energy. The second line includes
electron-electron repulsion and electron-nucleus attraction. For nuclei of
general charges, the corresponding factors of \(Z_a\) are restored; the
lecture keeps the charges schematic because the separation of time scales is
the issue.

There are two characteristic motions. The nuclei are heavy and slow; the
electrons are light and fast enough to look like a blur. In the first
approximation, fix the nuclear positions and group all electron-dependent
terms into an electronic Hamiltonian,
\begin{equation}
H_{\mathrm{el}}(R_1,R_2).
\end{equation}
For each fixed nuclear configuration, solve
\begin{equation}
H_{\mathrm{el}}(R_1,R_2)
\psi_0(r_1,\ldots,r_N;R_1,R_2)
=
E_{\mathrm{el}}(R_1,R_2)\psi_0.
\label{eq:electronic-ground-state}
\end{equation}
We retain the lowest electronic eigenvalue. The electrons have not been
ignored; all their rapid motion has been summarized by the function
\(E_{\mathrm{el}}(R_1,R_2)\).

The logic is adiabatic. During one rapid electronic adjustment, the nuclei
have moved so little that their coordinates can be treated as parameters.
After the electronic eigenproblem is solved, those same coordinates are
allowed to move again. At each slowly changing configuration the electron
cloud follows its instantaneous ground state. A more accurate treatment
could retain excited electronic surfaces and transitions between them, but
the leading surface already exhibits the renormalization idea.

\begin{figure}[htbp]
\centering
\includegraphics[width=0.72\textwidth]{lecture_01_figure_02.png}
\caption{The molecular Hamiltonian remains on the board as the electronic
ground-state problem is introduced. The typed form in
Eq.~\eqref{eq:molecular-hamiltonian} records the terms obscured in this frame.}
\lectureframe{00:16:19}
\end{figure}

\begin{classroomqa}
\audiencequestion{Are the nuclear momenta simply being set to
zero?\lecturetimestamp{00:15:39}}
\lecturerresponse{Only temporarily. In the leading approximation the nuclei
are treated as nailed down because they move much more slowly than the
electrons. Their momenta will be restored when we solve the effective nuclear
problem, so this approximation is not a violation of the uncertainty
principle.}
\end{classroomqa}

The electronic ground-state energy is now a contribution to the potential
felt by the nuclei:
\begin{equation}
V_{\mathrm{eff}}(R_1,R_2)
=\frac{e^2}{R_{12}}+E_{\mathrm{el}}(R_1,R_2).
\label{eq:molecular-effective-potential}
\end{equation}
Rotational and translational invariance reduce it, for two isolated atoms, to
a function of their separation \(R_{12}\). At very short distance the nuclear
Coulomb repulsion dominates. At very large distance there is little
interaction. Between these limits the electronic cloud can produce an
attractive region and a minimum. The nuclear wavefunction then sits near that
minimum and describes a bound molecule.

Notice the order of operations. We did not guess an empirical attraction and
discard the electrons. We first solved their quantum problem in the fixed
nuclear background. Only then did we erase their explicit coordinates and
retain the resulting eigenvalue. The short-distance physics survives in the
shape, depth, and curvature of \(V_{\mathrm{eff}}\). The location of its
minimum sets the equilibrium separation, its curvature controls small
vibrations, and its asymptotic form controls whether a bound state exists.
One effective function carries all of that electronic information into the
slower problem.

Thus a complicated many-electron Hamiltonian has become a two-body nuclear
Hamiltonian with a new potential. We eliminated the fast degrees of freedom,
but retained their complete low-energy effect. The potential has been
renormalized. This familiar Born--Oppenheimer reasoning is not usually
introduced under that name, yet it contains precisely the idea we need:
\begin{equation}
\text{fast microscopic dynamics}
\quad\longrightarrow\quad
\text{effective parameters for slow dynamics}.
\end{equation}

\section{Natural Units and the Dimension of a Field}

In quantum field theory, the modes to be eliminated may be described
equivalently as short-distance, short-wavelength, high-frequency, or
high-energy modes. Before drawing diagrams, let us
fix the dimensional bookkeeping.

Ordinary mechanics uses independent units of length, time, and mass. In
particle physics we set
\begin{equation}
\hbar=c=1.
\end{equation}
Then energy and mass are equivalent by \(E=mc^2\), length and time are
equivalent because \(c=1\), and \(\hbar=1\) makes momentum the inverse of
length. Only one dimension remains:
\begin{equation}
[M]=[E]=[p]=L^{-1},
\qquad
[x]=[t]=L.
\label{eq:natural-units}
\end{equation}

This does not erase dimensions. It chooses conversion factors so that a time
can be quoted as a length and a mass as an inverse length. Once one unit is
specified, every other one follows. This is why a cutoff may be described
either as a short distance \(\delta\) or as a high energy
\(\Lambda=1/\delta\): they are the same scale in two languages.

Take one real scalar field \(\phi\), with a schematically normalized
Lagrangian density
\begin{equation}
\mathcal L=(\partial_\mu\phi)(\partial^\mu\phi)-V(\phi),
\qquad
S=\int d^4x\,\mathcal L.
\label{eq:scalar-action}
\end{equation}

\begin{figure}[htbp]
\centering
\includegraphics[width=0.78\textwidth]{lecture_01_figure_03.png}
\caption{The local scalar Lagrangian density and its spacetime integral into
the action.}
\lectureframe{00:26:08}
\end{figure}

The Lagrangian density is local: it describes one spacetime point. Classical
field equations follow by making the action stationary. In quantum theory the
same action appears in the path-integral phase \(\exp(iS)\). Either viewpoint
makes the dimension of \(S\) decisive. With \(\hbar\) restored, the phase is
\(\exp(iS/\hbar)\); after setting \(\hbar=1\), \(S\) itself must be
dimensionless.

\begin{classroomqa}
\audiencequestion{Does the action have units of
\(\hbar\)?\lecturetimestamp{00:26:13}}
\lecturerresponse{Yes. Since \(\hbar\) has been set equal to one, the action
is dimensionless in these units. The four-dimensional measure has dimension
\(L^4\), so the Lagrangian density must have dimension \(L^{-4}\).}
\end{classroomqa}

The kinetic term now fixes the dimension of the field:
\begin{align}
[S]&=1,
&
[\mathcal L]&=L^{-4},
&
[\partial_\mu]&=L^{-1},
\nonumber\\
[(\partial\phi)^2]&=L^{-2}[\phi]^2=L^{-4}
&\Longrightarrow\qquad
[\phi]&=L^{-1}.
\label{eq:scalar-dimension}
\end{align}
This small calculation will determine the ultraviolet size of every diagram
that follows.

Expand the potential as
\begin{equation}
V(\phi)=\frac{m^2}{2}\phi^2+g\phi^3+\lambda\phi^4+\cdots.
\label{eq:scalar-potential}
\end{equation}
Every term must have dimension \(L^{-4}\). Since
\([\phi^n]=L^{-n}\),
\begin{equation}
[m^2]=L^{-2},
\qquad
[g]=L^{-1},
\qquad
[\lambda]=1.
\label{eq:coupling-dimensions}
\end{equation}
The quartic coupling is dimensionless. Such couplings are especially
important in renormalization because their scale dependence cannot be read
off from a simple power of length.

\begin{classroomqa}
\audiencequestion{Does \(g\) have units of mass?\lecturetimestamp{00:29:52}}
\lecturerresponse{Yes. The factor \(\phi^3\) has dimension \(L^{-3}\), while
the potential density must have dimension \(L^{-4}\). Therefore
\([g]=L^{-1}\), which is a mass in natural units.}
\end{classroomqa}

\section{Vertices and the Short-Distance Propagator}

The Feynman rules are built from vertices and propagators. The term
\(g\phi^3\) gives a three-line vertex carrying a factor proportional to
\(g\); \(\lambda\phi^4\) gives a four-line vertex carrying a factor
proportional to \(\lambda\). The quadratic term gives a local two-point
insertion. It may be pictured as a particle entering, being absorbed, and
being re-emitted, with a coefficient proportional to \(m^2/2\). Numerical
normalizations, signs, factors of \(i\), and symmetry factors are deliberately
suppressed because the lecture is extracting scale dependence rather than
computing a cross section.

These drawings have an operational reading. A three-point vertex can absorb
two quanta and emit one, or represent any crossing-related assignment of its
three lines. A four-point vertex can describe two particles entering and two
leaving. The quadratic cross is still a legitimate two-line vertex even
though it is normally incorporated into the free propagator. We display it
because every unresolved two-point process will shortly be recognized by
comparison with it.

A propagator joins two spacetime points. Its line is not a claim that a
quantum particle followed a straight classical path. It represents the
amplitude to create a field excitation at \(x\) and remove or detect it at
\(y\):
\begin{equation}
\Delta(x-y)=\langle0|\phi(y)\phi(x)|0\rangle.
\label{eq:position-propagator}
\end{equation}
Read the expression from the right. Acting on the vacuum with \(\phi(x)\)
creates a one-particle component at \(x\). Acting with \(\phi(y)\) on the bra
removes the component at \(y\). Their vacuum matrix element is the overlap
amplitude between creation at the first point and detection at the second.
The symbols \(x\) and \(y\) denote whole four-vectors, not merely two
Cartesian coordinates.

States carry no physical dimension, while each field contributes \(L^{-1}\).
Therefore
\begin{equation}
[\Delta]=L^{-2}.
\end{equation}
For a massless field there is no intrinsic length scale, so dimensional
analysis forces the short-distance behavior
\begin{equation}
\Delta(x-y)\sim\frac{1}{(x-y)^2}.
\label{eq:massless-propagator}
\end{equation}
The denominator is the Lorentz-invariant squared separation, with the usual
prescription needed at its singularities.\editorialnote{For power counting it
is convenient to Wick rotate to Euclidean separation, where the magnitude is
unambiguously proportional to \(1/|x-y|^2\). This does not change the
ultraviolet powers used below.}

As \(y\) approaches \(x\), the propagator diverges. This singularity of
propagation over a vanishing separation is the source from which ultraviolet
divergences in more elaborate diagrams are built.

\section{How a Scalar Mass Is Renormalized}

For a scalar field the dynamical parameter is \(m^2\), as the relativistic
dispersion relation already suggests:
\begin{equation}
E=\sqrt{\mathbf p^2+m^2}.
\end{equation}
Let \(\delta\) be the shortest distance our effective description resolves.
We shall not distinguish processes whose vertices are separated by less than
\(\delta\). Any such unresolved process with one incoming and one outgoing
scalar has the same observable two-point shape as a mass insertion and must
therefore contribute to the effective mass.

Think of an accelerator as a microscope. If its finest resolution is
\(\delta\), an absorption at one point followed by an emission a distance
much smaller than \(\delta\) away cannot be distinguished from a truly local
event. We deliberately blur the vertex over that scale. Renormalization is
the rule that transfers all such sub-resolution alternatives into the
coefficient of the corresponding local operator.

\subsection{The first quartic correction}

Begin with a single \(\phi^4\) vertex. Two legs are the external particle;
the other two join into a short loop emitted and reabsorbed at nearly the same
point. Smearing the coincident propagator over the cutoff distance gives
\begin{equation}
\Delta(0)\ \longrightarrow\ \Delta(\delta)\sim\frac{1}{\delta^2},
\end{equation}
and hence
\begin{equation}
\delta m^2_{(1)}\sim\frac{\lambda}{\delta^2}.
\label{eq:one-loop-scalar-mass}
\end{equation}
If \(m_0^2\) denotes the input or bare parameter, the low-resolution
two-point coefficient has the schematic form
\begin{equation}
\frac{m_{\mathrm{eff}}^2}{2}
\sim
\frac{m_0^2}{2}
+c_1\frac{\lambda}{\delta^2},
\label{eq:first-renormalized-mass}
\end{equation}
where \(c_1\) contains the ignored numerical and convention-dependent
factors. The laboratory mass measured in experiments that do not resolve
distances below \(\delta\) is this renormalized quantity, not \(m_0\) by
itself.

There are therefore two indistinguishable ways for the observed two-point
event to occur. The original quadratic vertex can absorb and re-emit the
particle, or the quartic interaction can do so while briefly emitting and
reabsorbing a virtual scalar. Quantum mechanics adds the amplitudes. Once the
loop lies below the resolution scale, its separate history is gone from the
effective description, but its contribution to the measured coefficient is
not.

\subsection{The next quartic correction}

Now connect two quartic vertices so that one external particle enters at
\(x\), one leaves from \(x+\Delta\), and three internal propagators join the
vertices. There are two powers of \(\lambda\), while the propagators give
\(\Delta^{-6}\). Quantum amplitudes must be summed over every unresolved
place from which the outgoing particle can emerge, so the relative
four-vector is integrated:
\begin{equation}
\delta m^2_{(2)}
\sim
\lambda^2
\int_{|\Delta|\gtrsim\delta}
d^4\Delta\,\frac{1}{|\Delta|^6}.
\label{eq:two-loop-power-count}
\end{equation}

\begin{classroomqa}
\audiencequestion{Does the small mark at the vertex mean there are only three
particles there?\lecturetimestamp{00:46:35}}
\lecturerresponse{No. We are still using the quartic interaction. Four lines
meet at each vertex; in this two-point diagram one line is external and three
run between the two vertices. The second vertex again has four incident
lines.}
\end{classroomqa}

The measure supplies four powers of length and the three propagators supply
minus six. Since the ultraviolet cutoff is the only scale in the estimate,
\begin{equation}
\int_{|\Delta|\gtrsim\delta}
\frac{d^4\Delta}{|\Delta|^6}
\sim\frac{1}{\delta^2},
\qquad
\delta m^2_{(2)}
\sim\frac{\lambda^2}{\delta^2}.
\label{eq:two-loop-scalar-mass}
\end{equation}
Here \(\Delta^\mu\) is a four-component integration variable; one integrates
over its time and three spatial components. The lowercase \(\delta\) is the
fixed cutoff. The lecture initially uses the same spoken name for both and
later corrects the notation.

The same result is visible by writing the four-dimensional measure radially.
Up to an angular constant,
\begin{equation}
\int_{\delta}\frac{d^4\Delta}{|\Delta|^6}
\sim
\int_{\delta}\frac{R^3\,dR}{R^6}
=\int_{\delta}\frac{dR}{R^3}
\sim\frac{1}{\delta^2}.
\end{equation}
This is why the large-distance upper limit is unimportant for the ultraviolet
estimate: the integral is dominated by its lower endpoint.

Higher diagrams repeat the pattern. A diagram with more quartic vertices adds
a higher power of \(\lambda\), but dimensional analysis still permits a
quadratically divergent mass contribution:
\begin{equation}
m_{\mathrm{eff}}^2
=m_0^2
+\frac{1}{\delta^2}
\left(c_1\lambda+c_2\lambda^2+c_3\lambda^3+\cdots\right)
+\text{weaker terms}.
\label{eq:scalar-mass-series}
\end{equation}

\subsection{The cubic interaction and a logarithm}

Two cubic vertices contain two propagators between nearby points. Their
contribution scales as
\begin{equation}
\delta m_g^2
\sim
g^2\int_{|\Delta|\gtrsim\delta}
\frac{d^4\Delta}{|\Delta|^4}.
\label{eq:cubic-loop}
\end{equation}
Now the measure and denominator carry equal powers. The remaining dependence
is logarithmic rather than a power:
\begin{equation}
\delta m_g^2
\sim
c_g g^2\log\!\left(\frac{1}{\mu\delta}\right).
\label{eq:cubic-log}
\end{equation}
The reference scale \(\mu\) makes the logarithm dimensionless. This is the
fully dimensional version of the lecture's schematic \(g^2\log\delta\). A
logarithm varies much more slowly than any inverse power, so
\(1/\delta^2\) dominates as the cutoff becomes short.

Again the radial form makes the distinction transparent:
\begin{equation}
\int_{\delta}^{R_{\mathrm{IR}}}
\frac{d^4\Delta}{|\Delta|^4}
\sim
\int_{\delta}^{R_{\mathrm{IR}}}\frac{R^3\,dR}{R^4}
=\log\!\left(\frac{R_{\mathrm{IR}}}{\delta}\right).
\end{equation}
Equal powers in measure and denominator produce a logarithm. More powers in
the denominator produce an inverse power of the cutoff. The lecture is not
asking us to evaluate angular factors; it is teaching us to recognize the
ultraviolet answer on sight.

\begin{classroomqa}
\audiencequestion{To which base is the logarithm
taken?\lecturetimestamp{00:57:42}}
\lecturerresponse{Use the natural logarithm. Changing the base only multiplies
it by a numerical constant, and that constant is absorbed into the coefficient
whose factors of \(2\), \(\pi\), and symmetry numbers have not been
calculated here.}
\end{classroomqa}

\begin{classroomqa}
\audiencequestion{How can one take the logarithm of a length, and where do the
units of the mass correction come from?\lecturetimestamp{00:58:32}}
\lecturerresponse{A logarithm always compares two scales; properly it is the
logarithm of a dimensionless ratio such as \(1/(\mu\delta)\). The coupling
\(g\) has one power of mass, so \(g^2\) already has exactly the mass-squared
dimension required for the correction.}
\end{classroomqa}

\begin{classroomqa}
\audiencequestion{A Feynman diagram gives an amplitude. Is its squared
magnitude a probability?\lecturetimestamp{00:59:41}}
\lecturerresponse{Yes, when the amplitude is placed in a physical transition
calculation with the required state and phase-space factors. Here its more
immediate role is different: an unresolved two-point amplitude has the same
local form as the mass term and therefore shifts the effective mass
parameter.}
\end{classroomqa}

\section{Every Parameter Becomes Effective}

The infinite family of mass diagrams does not disappear merely because
\(\lambda\) is small. Perturbation theory orders
them by powers of the coupling, so successive terms may become smaller, but
there are still infinitely many. The same short-distance elimination also
renormalizes the interactions themselves.

\begin{classroomqa}
\audiencequestion{Even for mass renormalization, are there not infinitely
many diagrams?\lecturetimestamp{01:00:35}}
\lecturerresponse{There are. If \(\lambda\) is sufficiently small, the
perturbative series may be useful because higher powers are successively
suppressed. Nevertheless, the renormalized parameter is defined by the whole
series, not by the first drawing alone.}
\end{classroomqa}

Operationally, a local \(\lambda\phi^4\) interaction is a
two-particle-in, two-particle-out amplitude. If a more complicated four-point
process occurs entirely inside an unresolved region, it is indistinguishable
at low resolution from the original local vertex. Such diagrams are
therefore lumped into an effective \(\lambda\). Their leading scale dependence
is typically logarithmic in four-dimensional \(\phi^4\) theory. The lecture
does not calculate that beta function; it uses the example to make the
general statement:
\begin{equation}
\text{parameters measured at low energy}
\ne
\text{parameters entered at the microscopic cutoff}.
\end{equation}

The four-point example matters because it removes a possible
misunderstanding. Renormalization is not a special repair applied only to a
mass divergence. Every operator allowed in the Lagrangian can be reproduced
by unresolved diagrams with the same external legs and symmetries. A
two-point process renormalizes a quadratic coefficient, a four-point process
renormalizes a quartic coefficient, and vacuum-to-vacuum processes will
renormalize the constant term. The low-energy Lagrangian is the compact
ledger in which all of those short-distance processes are recorded.

\begin{classroomqa}
\audiencequestion{Since \(\Delta\) has units of length, is a logarithm of it
not dimensionally strange? And did the board interchange uppercase and
lowercase delta?\lecturetimestamp{01:01:34}}
\lecturerresponse{The logarithm really means a ratio to another scale and is
therefore dimensionless. Yes, the notation was briefly interchanged:
\(\Delta^\mu\) is the integration variable, while \(\delta\) is the lower
distance cutoff that remains in the answer.}
\end{classroomqa}

This is the field-theory counterpart of the molecule. We solved for the
unseen fast electrons and put their result into a potential. Here we sum
unseen short-distance processes and put their result into masses, couplings,
and every other allowed local coefficient. Everything is shifted.

\begin{classroomqa}
\audiencequestion{Could you repeat whether the measured quantity includes
scales smaller than the experiment can see?\lecturetimestamp{01:05:13}}
\lecturerresponse{It does. Their individual fluctuations are not resolved,
but their summed effect is present. All diagrams involving shorter scales are
collected into the single effective coefficient that a low-energy experiment
calls, for example, the physical mass squared.}
\end{classroomqa}

\begin{classroomqa}
\audiencequestion{Must the input term \(m_0^2\) then cancel the other
contributions?\lecturetimestamp{01:06:00}}
\lecturerresponse{If the measured scalar mass is far smaller than the
individual ultraviolet contributions, yes. The required cancellation is the
fine-tuning problem to which we now turn.}
\end{classroomqa}

\section{The Higgs Hierarchy Problem}

Imagine that the cutoff is the Planck length, so that the effective field
theory is being asked to summarize fluctuations from laboratory scales down
to the scale where quantum gravity must intervene.
In particle-physics units,
\begin{equation}
\delta\sim M_{\mathrm{Pl}}^{-1}\sim10^{-19}\,\mathrm{GeV}^{-1},
\qquad
\frac{1}{\delta^2}\sim10^{38}\,\mathrm{GeV}^2.
\label{eq:planck-cutoff}
\end{equation}
The precise use of the reduced or unreduced Planck scale is irrelevant to the
enormous comparison.

The distance hierarchy is equally striking. The Planck length is about
nineteen orders of magnitude below a proton radius and roughly sixteen or
seventeen orders below the distances directly probed by the accelerators
being discussed. Squaring that separation is what turns a large hierarchy of
lengths into the roughly thirty-four-place cancellation in a scalar
mass-squared parameter.

For the Higgs field, a scalar, the desired weak scale is of order a few
hundred GeV, not \(10^{19}\) GeV. Yet the terms in
Eq.~\eqref{eq:scalar-mass-series} can naturally be of Planck-scale size.
Their signs have not been tracked: some may be positive and some negative,
and the bare mass squared may also have either sign. But to leave an
electroweak-scale remainder, those huge pieces must cancel to roughly
thirty-four decimal places. Susskind first says seventeen digits, then
corrects himself because the comparison is between squared masses. It is the
squared ratio that measures this version of the tuning.

\begin{classroomqa}
\audiencequestion{Is \(\lambda\) smaller or larger than
one?\lecturetimestamp{01:09:50}}
\lecturerresponse{It is assumed to be modestly smaller than one. Perturbation
theory would become intractable if successive powers grew, and a pure scalar
theory also cannot be extrapolated consistently with an arbitrarily large
quartic coupling. But \(\lambda\) is not remotely as small as
\(10^{-38}\); the lecture takes a percent-scale number for orientation.}
\end{classroomqa}

\begin{classroomqa}
\audiencequestion{Even if \(\lambda<1\), do not many terms remain important
because each carries the large cutoff factor?\lecturetimestamp{01:10:43}}
\lecturerresponse{Yes. Higher powers eventually suppress the series, but a
large collection of terms can contribute before that happens. They and the
bare term would all have to conspire to leave the extraordinarily small
physical answer.}
\end{classroomqa}

Recall why the Higgs mass parameter is known to be weak-scale. In the
lecture's convention, write the symmetry-breaking potential as
\begin{equation}
V(\phi)=-\frac{\mu^2}{2}\phi^2+\frac{\lambda}{4}\phi^4.
\label{eq:higgs-potential}
\end{equation}
The negative quadratic coefficient makes the origin unstable, while the
positive quartic term stabilizes the potential. Differentiating,
\begin{equation}
\frac{dV}{d\phi}=-\mu^2\phi+\lambda\phi^3
=\phi(-\mu^2+\lambda\phi^2),
\end{equation}
gives the nonzero minimum
\begin{equation}
\phi_{\min}^2=f^2=\frac{\mu^2}{\lambda}.
\label{eq:higgs-minimum}
\end{equation}
The board writes the quartic schematically as \(\lambda\phi^4\) but quotes
Eq.~\eqref{eq:higgs-minimum}; inserting the conventional factor \(1/4\)
makes those two lines consistent without changing the argument.

The sign of \(-\mu^2\) is important but does not rescue naturalness. It is the
negative curvature at the origin that drives spontaneous symmetry breaking.
Renormalization shifts that curvature by large positive and negative pieces.
The observed broken vacuum requires their sum not merely to be small, but to
have the particular small negative value that places the minimum at the weak
scale.

The \(W\)- and \(Z\)-boson masses determine \(f\) to be of order
\(200\,\mathrm{GeV}\) in the rough normalization used in the lecture. If
\(\lambda\) is neither enormous nor fantastically tiny, \(\mu\) must be of
the same general order. The observed weak-scale coefficient is therefore the
residue of contributions that may each be tens of orders of magnitude
larger.

This is called the Higgs fine-tuning problem, the hierarchy problem, or the
gauge-hierarchy problem. The hierarchy is the immense ratio between the
ordinary weak scale and a possible unification or Planck scale. A small
number by itself need not be paradoxical. What is disturbing here is a small
number assembled from large, apparently unrelated pieces that must cancel
with extreme precision. The problem was recognized in the late 1970s and
early 1980s, and the search for a mechanism behind that cancellation
motivates supersymmetry and several competing ideas.

The infinite perturbative series sharpens rather than removes the puzzle. A
percent-scale \(\lambda\) makes sufficiently high powers small, but many terms
can occur before that suppression wins. It is not enough for one loop to
cancel the bare mass accidentally. The effective weak scale must remain
stable after the rest of the short-distance processes are included. A
structural relation among the terms would be an explanation; an unrelated
digit-by-digit cancellation is only a tuning.

\section{Why a Fermion Mass Behaves Better}

The other Standard Model particles do not require independent tunings of the
same kind. For a Dirac fermion, the mass term is
\begin{equation}
\mathcal L_m=-m\bar\psi\psi
=-m\left(\bar\psi_L\psi_R+\bar\psi_R\psi_L\right).
\label{eq:fermion-mass}
\end{equation}
It contains the mass, not its square, and it couples left and right
chiralities. A local mass insertion therefore changes \(L\) into \(R\), or
\(R\) into \(L\).

\begin{figure}[htbp]
\centering
\includegraphics[width=0.76\textwidth]{lecture_01_figure_04.png}
\caption{The fermion mass term and the board's chirality-changing line
sketches. The diagrams are reconstructed algebraically below because part of
the chalk work is occluded.}
\lectureframe{01:23:36}
\end{figure}

The spoken lecture often says handedness or helicity. The exact protection is
most cleanly stated in terms of chirality. For an ultrarelativistic particle
the helicity picture supplies the intended intuition, but chirality is the
quantity carried by the left- and right-handed fields in
Eq.~\eqref{eq:fermion-mass}.

Gauge interactions preserve chirality:
\begin{equation}
L\xrightarrow{\ \gamma\ }L,
\qquad
R\xrightarrow{\ \gamma\ }R.
\label{eq:gauge-preserves-chirality}
\end{equation}
A Yukawa interaction with a scalar connects opposite chiralities:
\begin{equation}
L\xleftrightarrow{\ \phi\ }R.
\label{eq:scalar-flips-chirality}
\end{equation}
Suppose first that the electron starts with \(m_0=0\). Emitting and
reabsorbing any number of photons cannot produce a mass, because every gauge
vertex preserves chirality. A sequence that starts left remains left. Thus
pure quantum electrodynamics consistently permits a massless electron: its
electromagnetic self-energy cannot create an additive mass from zero.

In the line picture, start with a left-chiral electron. At the first photon
vertex it is still left-chiral; after propagation and reabsorption of the
photon it is still left-chiral. Repeating the construction only inserts more
chirality-preserving vertices. The external two-point function can be
renormalized, but it cannot acquire the left-right structure that defines a
mass.

Could scalar emission do it? One scalar vertex changes \(L\) to \(R\), but
an emitted scalar must be reabsorbed. A two-point self-energy graph therefore
contains an even number of scalar vertices. The second flip reverses the
first, so the complete graph again preserves the external chirality. No
matter how complicated the internal graph becomes, emission and reabsorption
pair the vertices and prevent an additive \(L\)-to-\(R\) term from appearing
when the starting mass is zero.

This is a parity argument on the number of flips. One scalar vertex flips
chirality once; closing every internal scalar line requires a second endpoint
and hence a second flip. Four scalar vertices flip four times, and so on.
An even number returns the fermion to its original chirality. The conclusion
is not restricted to the simplest loop drawn on the board.

A nonzero mass correction becomes possible only after the diagram already
contains a mass insertion. For example, photon emission and reabsorption can
surround one \(m_0\bar\psi\psi\) insertion. The gauge vertices contribute two
powers of the electric charge, but the sole chirality flip contributes the
original mass:
\begin{equation}
m_e
=m_0+c_1e^2m_0+c_2e^4m_0+\cdots
=m_0\left(1+c_1e^2+c_2e^4+\cdots\right).
\label{eq:multiplicative-fermion-mass}
\end{equation}
If \(m_0=0\), the result remains zero. If \(m_0\) is small, every correction
is proportional to that small number. We may still ask why the electron is
light relative to a fundamental scale, but we are no longer subtracting
enormous unrelated numbers to obtain it. In this sense a small fermion mass
is technically natural.

The mass insertion makes the allowed sequence explicit. Begin with \(L\).
Photon emission preserves \(L\); the insertion \(m_0\) changes \(L\) to
\(R\); photon reabsorption preserves \(R\). The diagram now has the external
chirality of a mass term, but its value necessarily contains the factor
\(m_0\). This is the crucial contrast with
Eq.~\eqref{eq:one-loop-scalar-mass}, whose scalar correction remains nonzero
even when the starting scalar mass is set to zero.

Gauge bosons have an analogous protection: gauge symmetry forbids an
independent mass term, so a gauge boson that begins massless remains massless
to all perturbative orders. In the Standard Model, the masses of fermions
and the weak gauge bosons arise in response to the Higgs expectation value.
Control the Higgs scale, and the remaining light masses do not require
separate copies of the same catastrophic tuning.

This identifies the question supersymmetry is meant to address: can a theory
prevent the Higgs scalar from receiving an enormous additive mass correction?
The answer is yes in suitably supersymmetric settings, because bosonic and
fermionic contributions are related rather than independent. The detailed
mechanism begins in the next lecture.

\section{The Other Fine-Tuning: Vacuum Energy}

There is one still more severe renormalization problem, but it becomes
observable only when gravity is included. In
nongravitational physics, adding a constant to every energy changes no
experiment. The electron mass, for example, is not an absolute energy; it is
the extra energy required to add an electron to the vacuum:
\begin{equation}
m_e=E(\text{vacuum}+\text{one electron})-E(\text{vacuum})
\end{equation}
in the electron rest frame. If the vacuum and every excited state shift
together, their differences remain unchanged.

Nevertheless, vacuum-to-vacuum diagrams do shift the vacuum energy. In
\(\lambda\phi^4\) theory, join the four legs of one quartic vertex into two
short loops. Each coincident propagator supplies \(1/\delta^2\), so
\begin{equation}
\delta\rho_{\mathrm{vac}}
\sim
c_{\mathrm{vac}}\frac{\lambda}{\delta^4}.
\label{eq:vacuum-energy-divergence}
\end{equation}
More complicated vacuum diagrams carry different coupling factors and
coefficients, but dimensional analysis typically leaves the same quartic
cutoff sensitivity for an energy density.

The absence of external lines is exactly what identifies the correction.
Two-point graphs compare a one-particle state with the vacuum and alter a
mass; four-point graphs alter scattering; a graph with no incoming or
outgoing particles changes the vacuum-to-vacuum amplitude itself. One
quartic vertex supplies \(\lambda\), and its two short loops supply
\(\delta^{-2}\) apiece. Dimensional analysis then requires the
\(\delta^{-4}\) energy-density scale.

Without gravity one may choose the additive zero of energy and ask who cares.
General relativity removes that freedom: stress-energy sources spacetime
curvature, so a uniform vacuum energy makes empty space gravitate. It appears
as a cosmological constant and is connected observationally with dark energy.
A tiny measured vacuum energy must then survive contributions of the form
\(\delta^{-4}\).

This is why the question changes when gravity is turned on. In ordinary
quantum mechanics, shifting every level by the same constant leaves all
transition frequencies untouched. In gravity, the constant is not inert:
the vacuum stress tensor is proportional to the metric and curves spacetime.
What had been an arbitrary zero of energy becomes a source in Einstein's
equation.

Thus there are at least two extreme naturalness puzzles. The lecture quotes
roughly 123 decimal places for the vacuum-energy tuning and roughly 30--35
for the Higgs mass tuning. The exact counts depend on the ultraviolet scale
and conventions, but the contrast is not in doubt. Supersymmetry is a
potential response to the Higgs problem discussed here; once supersymmetry is
broken, it does not by itself explain the observed cosmological
constant.\editorialnote{This is also why an exactly supersymmetric
boson--fermion cancellation is not a phenomenological solution to the
observed vacuum energy: supersymmetry must be broken in our world.}

\section{Closing Classroom Questions}

\begin{classroomqa}
\audiencequestion{The Higgs correction depends on the cutoff \(\delta\), but
the cutoff seems arbitrary. Could it be the Planck length, ten times the
Planck length, or something else? Does that make the fine-tuning
arbitrary?\lecturetimestamp{01:33:50}}
\lecturerresponse{The numerical contribution depends on the shortest scale to
which the effective theory is trusted. That dependence is physical
bookkeeping, not an escape. Modes down to any scale where the Standard Model
still makes sense are modes whose effects must be included. Choose
\(\delta\) at the distance where new physics is expected---perhaps the Planck
scale, perhaps a unification scale near \(10^{16}\,\mathrm{GeV}\). The
contribution accumulated down to that scale is already a lower bound on the
problem. Extending the same theory to still shorter distances only makes the
quadratic sensitivity larger. We are obligated to explain what cancels the
contribution from every scale between the laboratory and that ultraviolet
boundary.}
\end{classroomqa}

\begin{classroomqa}
\audiencequestion{What defines a scalar particle, and what is an
example?\lecturetimestamp{01:37:26}}
\lecturerresponse{The Higgs is the intended example. A scalar does not
transform into different components under a spatial rotation; equivalently,
its particle quanta have spin zero. A photon is described by a vector field,
and the gravitational field by a tensor. At the time of this 2010 lecture, no
elementary spin-zero particle had yet been detected.\editorialnote{The Higgs
boson was discovered at CERN in 2012 and its measured properties are
consistent with spin zero.}}
\end{classroomqa}

\begin{classroomqa}
\audiencequestion{Is the \(Z\) particle also a gauge
boson?\lecturetimestamp{01:38:31}}
\lecturerresponse{Yes. The \(Z\), \(W^+\), and \(W^-\) are spin-one gauge
bosons, as is the photon. Their field equations and interactions are
non-Abelian or electroweak generalizations of the Maxwell-like gauge
description.}
\end{classroomqa}

\begin{classroomqa}
\audiencequestion{Why does the bosonic ground-state contribution not cancel
the fermionic one in the vacuum energy?\lecturetimestamp{01:39:10}}
\lecturerresponse{Fermion loops and corresponding boson loops can indeed
enter with opposite signs. Cancellation is not automatic, however: the
particle masses, multiplicities, and couplings must match. Even a small
mismatch leaves a large cutoff-sensitive residue, and the infinite families
of diagrams must match order by order. Exact supersymmetry enforces just such
boson--fermion relations. Generic nonsupersymmetric particle content does
not, and broken supersymmetry leaves a residual vacuum-energy problem.}
\end{classroomqa}

\section{Perspective}

The route through the lecture is deliberate. A molecular calculation first
shows that eliminating fast variables is neither neglect nor magic: their
ground-state energy becomes a potential for the slow variables. Quantum
field theory repeats that operation across wavelengths. Dimensional analysis
then reveals which effective parameters are sensitive to the cutoff. Scalar
masses receive additive terms proportional to \(\delta^{-2}\); a small Higgs
scale therefore appears unnaturally delicate. Fermion masses are
multiplicatively renormalized because chirality protects the zero-mass limit,
and gauge symmetry similarly protects massless gauge bosons. Vacuum energy
is more sensitive still, scaling as \(\delta^{-4}\), and gravity makes its
absolute value observable.

Supersymmetry enters the course at exactly this point, not as ornamental new
mathematics but as an attempt to make the scalar sector behave more like the
protected sectors. The next task is to understand what symmetry can relate
bosons and fermions strongly enough to organize those cancellations.
